\begin{longtable}{ | p{5.5cm} | p{6.5cm} | p{1.3cm} | p{1.8cm} | }
\hline\rowcolor{lightgray}
\textbf{名称} & \textbf{描述} & \textbf{地址} & \textbf{访问} \\ \hline
\endhead
\multicolumn{4}{|l|}{\textbf{FIFO 读/写寄存器}} \\ \hline
\hyperref[regdesc:RMTCHNDATAREG]{RMT\_CH0DATA\_REG} & 通道 0 通过 APB FIFO 进行读写操作时用到的数据寄存器 & 0x0000 & HRO \\ \hline
\hyperref[regdesc:RMTCHNDATAREG]{RMT\_CH1DATA\_REG} & 通道 1 通过 APB FIFO 进行读写操作时用到的数据寄存器 & 0x0004 & HRO \\ \hline
\hyperref[regdesc:RMTCHNDATAREG]{RMT\_CH2DATA\_REG} & 通道 2 通过 APB FIFO 进行读写操作时用到的数据寄存器 & 0x0008 & HRO \\ \hline
\hyperref[regdesc:RMTCHNDATAREG]{RMT\_CH3DATA\_REG} & 通道 3 通过 APB FIFO 进行读写操作时用到的数据寄存器 & 0x000C & HRO \\ \hline
\hyperref[regdesc:RMTCHMDATAREG]{RMT\_CH4DATA\_REG} & 通道 4 通过 APB FIFO 进行读写操作时用到的数据寄存器 & 0x0010 & HRO \\ \hline
\hyperref[regdesc:RMTCHMDATAREG]{RMT\_CH5DATA\_REG} & 通道 5 通过 APB FIFO 进行读写操作时用到的数据寄存器 & 0x0014 & HRO \\ \hline
\hyperref[regdesc:RMTCHMDATAREG]{RMT\_CH6DATA\_REG} & 通道 6 通过 APB FIFO 进行读写操作时用到的数据寄存器 & 0x0018 & HRO \\ \hline
\hyperref[regdesc:RMTCHMDATAREG]{RMT\_CH7DATA\_REG} & 通道 7 通过 APB FIFO 进行读写操作时用到的数据寄存器 & 0x001C & HRO \\ \hline
\multicolumn{4}{|l|}{\textbf{配置寄存器}} \\ \hline
\hyperref[regdesc:RMTCHNCONF0REG]{RMT\_CH0CONF0\_REG} & 通道 0 的配置寄存器 0 & 0x0020 & varies \\ \hline
\hyperref[regdesc:RMTCHNCONF0REG]{RMT\_CH1CONF0\_REG} & 通道 1 的配置寄存器 0 & 0x0024 & varies \\ \hline
\hyperref[regdesc:RMTCHNCONF0REG]{RMT\_CH2CONF0\_REG} & 通道 2 的配置寄存器 0 & 0x0028 & varies \\ \hline
\hyperref[regdesc:RMTCHNCONF0REG]{RMT\_CH3CONF0\_REG} & 通道 3 的配置寄存器 0 & 0x002C & varies \\ \hline
\hyperref[regdesc:RMTCHMCONF0REG]{RMT\_CH4CONF0\_REG} & 通道 4 的配置寄存器 0 & 0x0030 & R/W \\ \hline
\hyperref[regdesc:RMTCHMCONF1REG]{RMT\_CH4CONF1\_REG} & 通道 4 的配置寄存器 1 & 0x0034 & varies \\ \hline
\hyperref[regdesc:RMTCHMCONF0REG]{RMT\_CH5CONF0\_REG} & 通道 5 的配置寄存器 0 & 0x0038 & R/W \\ \hline
\hyperref[regdesc:RMTCHMCONF1REG]{RMT\_CH5CONF1\_REG} & 通道 5 的配置寄存器 1 & 0x003C & varies \\ \hline
\hyperref[regdesc:RMTCHMCONF0REG]{RMT\_CH6CONF0\_REG} & 通道 6 的配置寄存器 0 & 0x0040 & R/W \\ \hline
\hyperref[regdesc:RMTCHMCONF1REG]{RMT\_CH6CONF1\_REG} & 通道 6 的配置寄存器 1 & 0x0044 & varies \\ \hline
\hyperref[regdesc:RMTCHMCONF0REG]{RMT\_CH7CONF0\_REG} & 通道 7 的配置寄存器 0 & 0x0048 & R/W \\ \hline
\hyperref[regdesc:RMTCHMCONF1REG]{RMT\_CH7CONF1\_REG} & 通道 7 的配置寄存器 1 & 0x004C & varies \\ \hline
\hyperref[regdesc:RMTCHMRXCARRIERRMREG]{RMT\_CH4\_RX\_CARRIER\_RM\_REG} & 通道 4 的载波滤除寄存器 & 0x0090 & R/W \\ \hline
\hyperref[regdesc:RMTCHMRXCARRIERRMREG]{RMT\_CH5\_RX\_CARRIER\_RM\_REG} & 通道 5 的载波滤除寄存器 & 0x0094 & R/W \\ \hline
\hyperref[regdesc:RMTCHMRXCARRIERRMREG]{RMT\_CH6\_RX\_CARRIER\_RM\_REG} & 通道 6 的载波滤除寄存器 & 0x0098 & R/W \\ \hline
\hyperref[regdesc:RMTCHMRXCARRIERRMREG]{RMT\_CH7\_RX\_CARRIER\_RM\_REG} & 通道 7 的载波滤除寄存器 & 0x009C & R/W \\ \hline
\hyperref[regdesc:RMTSYSCONFREG]{RMT\_SYS\_CONF\_REG} & RMT APB 配置寄存器 & 0x00C0 & R/W \\ \hline
\hyperref[regdesc:RMTREFCNTRSTREG]{RMT\_REF\_CNT\_RST\_REG} & RMT 时钟分频器复位寄存器 & 0x00C8 & WT \\ \hline
\multicolumn{4}{|l|}{\textbf{状态寄存器}} \\ \hline
\hyperref[regdesc:RMTCHNSTATUSREG]{RMT\_CH0STATUS\_REG} & 通道 0 的状态寄存器 & 0x0050 & RO \\ \hline
\hyperref[regdesc:RMTCHNSTATUSREG]{RMT\_CH1STATUS\_REG} & 通道 1 的状态寄存器 & 0x0054 & RO \\ \hline
\hyperref[regdesc:RMTCHNSTATUSREG]{RMT\_CH2STATUS\_REG} & 通道 2 的状态寄存器 & 0x0058 & RO \\ \hline
\hyperref[regdesc:RMTCHNSTATUSREG]{RMT\_CH3STATUS\_REG} & 通道 3 的状态寄存器 & 0x005C & RO \\ \hline
\hyperref[regdesc:RMTCHMSTATUSREG]{RMT\_CH4STATUS\_REG} & 通道 4 的状态寄存器 & 0x0060 & RO \\ \hline
\hyperref[regdesc:RMTCHMSTATUSREG]{RMT\_CH5STATUS\_REG} & 通道 5 的状态寄存器 & 0x0064 & RO \\ \hline
\hyperref[regdesc:RMTCHMSTATUSREG]{RMT\_CH6STATUS\_REG} & 通道 6 的状态寄存器 & 0x0068 & RO \\ \hline
\hyperref[regdesc:RMTCHMSTATUSREG]{RMT\_CH7STATUS\_REG} & 通道 7 的状态寄存器 & 0x006C & RO \\ \hline
\multicolumn{4}{|l|}{\textbf{中断寄存器}} \\ \hline
\hyperref[regdesc:RMTINTRAWREG]{RMT\_INT\_RAW\_REG} & 原始中断状态寄存器 & 0x0070 & R/WTC/SS \\ \hline
\hyperref[regdesc:RMTINTSTREG]{RMT\_INT\_ST\_REG} & 屏蔽中断状态寄存器 & 0x0074 & RO \\ \hline
\hyperref[regdesc:RMTINTENAREG]{RMT\_INT\_ENA\_REG} &  中断使能寄存器 & 0x0078 & R/W \\ \hline
\hyperref[regdesc:RMTINTCLRREG]{RMT\_INT\_CLR\_REG} & 中断清零寄存器 & 0x007C & WT \\ \hline
\multicolumn{4}{|l|}{\textbf{载波占空比寄存器}} \\ \hline
\hyperref[regdesc:RMTCHNCARRIERDUTYREG]{RMT\_CH0CARRIER\_DUTY\_REG} & 通道 0 的占空比配置寄存器 & 0x0080 & R/W \\ \hline
\hyperref[regdesc:RMTCHNCARRIERDUTYREG]{RMT\_CH1CARRIER\_DUTY\_REG} & 通道 1 的占空比配置寄存器 & 0x0084 & R/W \\ \hline
\hyperref[regdesc:RMTCHNCARRIERDUTYREG]{RMT\_CH2CARRIER\_DUTY\_REG} & 通道 2 的占空比配置寄存器 & 0x0088 & R/W \\ \hline
\hyperref[regdesc:RMTCHNCARRIERDUTYREG]{RMT\_CH3CARRIER\_DUTY\_REG} & 通道 3 的占空比配置寄存器 & 0x008C & R/W \\ \hline
\multicolumn{4}{|l|}{\textbf{TX 事件配置寄存器}} \\ \hline
\hyperref[regdesc:RMTCHNTXLIMREG]{RMT\_CH0\_TX\_LIM\_REG} & 通道 0 的 TX 事件配置寄存器 & 0x00A0 & varies \\ \hline
\hyperref[regdesc:RMTCHNTXLIMREG]{RMT\_CH1\_TX\_LIM\_REG} & 通道 1 的 TX 事件配置寄存器 & 0x00A4 & varies \\ \hline
\hyperref[regdesc:RMTCHNTXLIMREG]{RMT\_CH2\_TX\_LIM\_REG} & 通道 2 的 TX 事件配置寄存器 & 0x00A8 & varies \\ \hline
\hyperref[regdesc:RMTCHNTXLIMREG]{RMT\_CH3\_TX\_LIM\_REG} & 通道 3 的 TX 事件配置寄存器 & 0x00AC & varies \\ \hline
\hyperref[regdesc:RMTTXSIMREG]{RMT\_TX\_SIM\_REG} & RMT TX 同步发送寄存器 & 0x00C4 & R/W \\ \hline
\multicolumn{4}{|l|}{\textbf{RX 事件配置寄存器}} \\ \hline
\hyperref[regdesc:RMTCHMRXLIMREG]{RMT\_CH4\_RX\_LIM\_REG} & 通道 4 RX 事件配置寄存器 & 0x00B0 & R/W \\ \hline
\hyperref[regdesc:RMTCHMRXLIMREG]{RMT\_CH5\_RX\_LIM\_REG} & 通道 5 RX 事件配置寄存器 & 0x00B4 & R/W \\ \hline
\hyperref[regdesc:RMTCHMRXLIMREG]{RMT\_CH6\_RX\_LIM\_REG} & 通道 6 RX 事件配置寄存器 & 0x00B8 & R/W \\ \hline
\hyperref[regdesc:RMTCHMRXLIMREG]{RMT\_CH7\_RX\_LIM\_REG} & 通道 7 RX 事件配置寄存器 & 0x00BC & R/W \\ \hline
\multicolumn{4}{|l|}{\textbf{版本寄存器}} \\ \hline
\hyperref[regdesc:RMTDATEREG]{RMT\_DATE\_REG} & 版本控制寄存器 & 0x00CC & R/W \\ \hline
\end{longtable}
